\section{Related work}
\label{sec:related}


\myparagraph{Structure-from-Motion and MVS.} 
SfM ~\cite{hartley2003multiple, crandall2012sfm, jiang2013global,cui2017hsfm, schonberger2016structure} builds 3D sparse maps by optimizing camera parameters using only RGB images. 
Traditional approaches consists of handcrafted pipelines to match keypoints ~\cite{harris1988combined, lowe2004distinctive, bay2006surf, rosten2006machine, barroso2019key}, form pixel correspondences across multiple images, determine geometric relations, and perform bundle adjustment to optimize camera locations.
Recently, SfM has incorporated learning-based approaches leveraging neural networks as components~\cite{yi2016lift, barath2023affineglue, chen2022aspanformer, lindenberger2023lightglue, sarlin2020superglue, sun2021loftr, tang2022quadtree, wang2023guiding}, but it remains brittle to outliers and challenging conditions.
Likewise, multi-view stereo (MVS) aims to reconstruct dense 3D surface through handcrafted pipelines based on triangulation from multiple viewpoints~\cite{furukawa2015multi, galliani2015massively, schonberger2016pixelwise}.
Learning-based approaches have also been incorporated for MVS to reconstruct 3D from few images~\cite{niemeyer2020differentiable,wiles2020synsin, yariv2020multiview, leroy2021volume, sitzmann2019srns,jang2021codenerf}, and more recently by adopting the RAFT architecture~\cite{cermvs} or using generic Transformers~\cite{cao2024mvsformer++}.
In contrast, the recent \duster{} framework~\cite{wang2024dust3r,leroy2024mast3r}, upon which we build our approach, departs from these methods altogether by defining a simple and unified neural framework to 3D reconstruction based on the paradigm of pointmap regression.
The idea is that, ideally, learning-based MVS should be capable of reconstructing 3D from a few images without strictly requiring the availability of camera parameters.

\myparagraph{Guiding 3D.} 
In this paper, we refer to \emph{guidance} as the notion that a method can  \emph{optionally} take in 3D priors (typically camera parameters or sparse depth) to help with some prediction task.
This can take different forms in the vast literature of 3D vision works.
COLMAP~\cite{schonberger2016pixelwise}, for instance, can take camera parameters for rectifying camera views; however, it is an optimization-based method, meaning that each scene has to be processed independently. %
Recently, UniDepth~\cite{piccinelli2024unidepth} is a monocular depth and intrinsic estimation method capable of optionally providing camera intrinsics as additional guide to generate depthmap, with an architecture specifically handcrafted for that purpose.
Also, depth completion models~\cite{guidenet,bpnet,yan2024tri,lrru} such as CompletionFormer~\cite{completionformer2023} take sparse depth input as guidance in order to inpaint the full depthmap, but are also typically capable of inferring depth from an image alone.
Overall, however, there is currently no learning-based model to the best of our knowledge that guides its output by taking any subset of sparse depth, camera intrinsics and camera extrinsics.


\myparagraph{Depth Completion.} 
Given sparse depths and a corresponding RGB image, the task is to densify  depth by propagating sparse measurements.
CNN-based approaches were adopted by ~\cite{ma2018self, ma2018sparse}, and ~\cite{cheng2018depth} proposed convolutional spatial propagation network (CSPN). 
Multi-branch networks, which process RGB and sparse depth separately, were used in ~\cite{hu2021penet, nazir2022semattnet, Qiu_2019_CVPR, guidenet, wvangansbeke2019sparse, zhang2018deepdepth, rho2022guideformer, yan2023rignet}.
CompletionFormer~\cite{completionformer2023} leverages both CNN and ViT architectures to capture both global and local features.
Most methods focus only on the specific depth completion task from a single image and partial depth as the input, often without incorporating camera intrinsics.
Using the pointmap representation, \ours{}, based on a ViT architecture, can complete the depth map from both RGB images and sparse point clouds, while taking relative pose or camera intrinsics as auxiliary input.


\myparagraph{Pose Estimation from Sparse views.} 
Learning-based approaches involve regressing relative pose, applying probabilistic methods, using energy-based or diffusion models~\cite{rockwell2022, zhang2022relpose, lin2024relpose++, balntas2017hpatches, wang2023posediffusion}.
RayDiffusion~\cite{raydiffusion} models camera rays instead of camera parameters, applying diffusion on rays, then recovers the camera pose using DLT~\cite{abdel1971direct}.
\duster{} discovers camera pose from pointmaps using PnP, while \ours{} predicts pointmaps in both camera coordinate frames, allowing the use of Procrustes alignment.

\myparagraph{\duster{} and its developments.} \duster{}~\cite{wang2024dust3r} does not require calibrated images for 3D reconstruction, can handle extreme viewpoint changes, and deals with both single and a few images within the same framework.
Its pointmap representation eases the training by not having to deal with camera parameters directly.
MASt3R~\cite{leroy2024mast3r} improves \duster{} by also predicting dense local features, trained with an additional matching loss.
Splatt3R~\cite{smart2024splatt3r} marries MASt3R with GS. %
Spann3R~\cite{wang2024spann3r} predicts pointmaps in global coordinate and  adopts an external spatial memory that tracks relevant 3D information.
MonST3R~\cite{zhang2024monst3r} further trains  \duster{} on dynamic scenes and generates time-varying dynamic point clouds, as well as camera intrinsics and poses from a dynamic video.
In this work, we take a different and complementary direction and instead improve the core capabilities of \duster{}, which could benefit all aforementioned works.

